%! Polynomial Equations, Inequalities, and Modeling — Unit 2, Lesson 2.7 — Group Activity
\documentclass[10pt]{article}
\usepackage{precalculus-article}
\usepackage{precalculus-boxes}
\newcommand{\TallMath}[1]{$\displaystyle #1\rule[-1.4em]{0pt}{3.2em}$}

\begin{document}

\pageheader{Unit 2, Lesson 2.7}{Group Activity}

\vspace{0.12in}

\begin{scenariobox}[The Packaging Desk]{plumlight}
\small
Your group runs the packaging desk. Today's stock is a flat sheet of card
measuring \textbf{16 cm by 24 cm}. A square of side $x$ is cut from each of the
four corners and the flaps are folded up to make an open box. Every tier works
this same sheet --- no tier draws a graph.

\begin{center}
\begin{tikzpicture}[x=0.20cm, y=0.20cm]
  \draw[very thick, plum] (0,0) rectangle (24,16);
  \draw[fill=lilacmid, draw=plum] (0,0) rectangle (3,3);
  \draw[fill=lilacmid, draw=plum] (21,0) rectangle (24,3);
  \draw[fill=lilacmid, draw=plum] (0,13) rectangle (3,16);
  \draw[fill=lilacmid, draw=plum] (21,13) rectangle (24,16);
  \node[below] at (12,-0.4) {\small $24$ cm};
  \node[left] at (-0.6,8) {\small $16$ cm};
  \node[right] at (3.4,1.5) {\small $x$};
\end{tikzpicture}
\end{center}
\end{scenariobox}

\vspace{0.1in}

\boxguard
\begin{tcolorbox}[colback=white, colframe=black!40,
    title=\textbf{Tier R --- Remediate},
    fonttitle=\bfseries, arc=2mm, left=3mm, right=3mm, top=2mm, bottom=2mm]

\begin{enumerate}[itemsep=8pt]
  \item Solve each equation. Each one is already set against $0$, or gets there
        in one step.

        \begin{enumerate}[label=(\alph*), itemsep=4pt]
          \item $(x-3)(x+5) = 0$ \quad \blank{3.4cm}
          \item $x(x-7) = 0$ \quad \blank{3.4cm}
          \item $x^2 - 6x = 0$

                \begin{work}
                  x^2 - 6x &= 0 \\
                  x(x-6) &= 0
                \end{work}

                Solutions: \blank{3.4cm}
        \end{enumerate}

  \item Let $p(x) = (x+3)(x-2)$. Its real zeros are $-3$ and $2$. Fill in the
        sign chart, then answer.

        \smallskip
        \renewcommand{\arraystretch}{1.5}
        \begin{center}
        \begin{tabular}{|l|c|c|c|}
        \hline
        \textbf{Interval} & $(-\infty,-3)$ & $(-3,2)$ & $(2,\infty)$ \\
        \hline
        \textbf{Test value} & \blank{1.2cm} & \blank{1.2cm} & \blank{1.2cm} \\
        \hline
        \textbf{Sign of $p$} & \blank{1.2cm} & \blank{1.2cm} & \blank{1.2cm} \\
        \hline
        \end{tabular}
        \end{center}
        \smallskip

        The solution of $p(x) > 0$ is \blank{5cm}

  \item The sheet is 16 cm by 24 cm and $x$ is the corner cut.

        \begin{enumerate}[label=(\alph*), itemsep=4pt]
          \item A cut cannot have zero or negative length, so $x >$
                \blank{1.2cm}
          \item The two cuts on the \textbf{16 cm} side must not meet, so
                $x <$ \blank{1.2cm}
          \item The domain of the model is therefore \blank{3.4cm}
        \end{enumerate}
\end{enumerate}
\end{tcolorbox}

\vspace{0.1in}

\boxguard
\begin{tcolorbox}[colback=white, colframe=black!40,
    title=\textbf{Tier A --- Approaching Proficiency},
    fonttitle=\bfseries, arc=2mm, left=3mm, right=3mm, top=2mm, bottom=2mm]

\begin{enumerate}[itemsep=8pt]
  \item Solve $x^3 = 2x^2 + 8x$.

        \begin{work}
          x^3 - 2x^2 - 8x &= 0 \\
          x\big(x^2 - 2x - 8\big) &= 0 \\
          x(x-4)(x+2) &= 0
        \end{work}

        Solutions: \blank{4.4cm}

        Which solution would you have lost by dividing both sides by $x$ at the
        start? \quad \blank{2cm}

  \item Solve $(x+2)(x-1)(x-3) \le 0$.

        \smallskip
        \renewcommand{\arraystretch}{1.5}
        \begin{center}
        \begin{tabular}{|l|c|c|c|c|}
        \hline
        \textbf{Interval} & $(-\infty,-2)$ & $(-2,1)$ & $(1,3)$ & $(3,\infty)$ \\
        \hline
        \textbf{Test value} & \blank{1.2cm} & \blank{1.2cm} & \blank{1.2cm} & \blank{1.2cm} \\
        \hline
        \textbf{Sign} & \blank{1.2cm} & \blank{1.2cm} & \blank{1.2cm} & \blank{1.2cm} \\
        \hline
        \end{tabular}
        \end{center}
        \smallskip

        Solution in interval notation: \blank{5cm}

        Are the zeros included this time? \quad \textbf{Yes} / \textbf{No}
        \quad Why? \blank{4.4cm}

  \item Build the volume model for the 16 cm by 24 cm sheet.

        \begin{enumerate}[label=(\alph*), itemsep=4pt]
          \item Height \blank{1.2cm}, width $16 - $ \blank{1.2cm},
                length $24 - $ \blank{1.2cm}
          \item $V(x) = $ \blank{6cm}
          \item A 3 cm cut: $V(3) = $ \blank{2.4cm}

                Write the answer as a sentence with units.

                \vspace{0.05in}
                \writeline
        \end{enumerate}
\end{enumerate}
\end{tcolorbox}

\vspace{0.1in}

\boxguard
\begin{tcolorbox}[colback=white, colframe=black!40,
    title=\textbf{Tier E --- Extension},
    fonttitle=\bfseries, arc=2mm, left=3mm, right=3mm, top=2mm, bottom=2mm]

\begin{enumerate}[itemsep=8pt]
  \item Write $V(x) = x(16-2x)(24-2x)$ in expanded form, then state the domain
        the context forces.

        \begin{work}
          V(x) &= x(16-2x)(24-2x) \\
               &= x\big(384 - 32x - 48x + 4x^2\big) \\
               &= x\big(4x^2 - 80x + 384\big) \\
               &= 4x^3 - 80x^2 + 384x
        \end{work}

        Three separate restrictions: $x > $ \blank{1cm}, \quad
        $x < $ \blank{1cm} (16 cm side), \quad $x < $ \blank{1cm} (24 cm side).

        Domain: \blank{3.4cm}

  \item Confirm that a 2 cm cut gives a box of volume $480$
        cm\textsuperscript{3}.

        \begin{work}
          V(2) &= 2(16-4)(24-4) \\
               &= 2(12)(20) \\
               &= 480
        \end{work}

  \item Which cuts give a box holding \textbf{at least} $480$
        cm\textsuperscript{3}? Solve $V(x) \ge 480$ on the domain.

        \begin{work}
          4x^3 - 80x^2 + 384x &\ge 480 \\
          x^3 - 20x^2 + 96x - 120 &\ge 0 \\
          (x-2)\big(x^2 - 18x + 60\big) &\ge 0
        \end{work}

        \begin{work}
          x^2 - 18x + 60 &= 0 \\
          x &= \frac{18 \pm \sqrt{324 - 240}}{2} \\
            &= \frac{18 \pm 2\sqrt{21}}{2} \\
            &= 9 \pm \sqrt{21}
        \end{work}

        One of those two roots must be thrown out. Which one, and on what
        grounds? \quad \blank{6cm}

        \smallskip
        \renewcommand{\arraystretch}{1.5}
        \begin{center}
        \begin{tabular}{|l|c|c|c|}
        \hline
        \textbf{Interval in the domain} & $(0,\,2)$ & $(2,\,4.42)$ & $(4.42,\,8)$ \\
        \hline
        \textbf{Test value} & \blank{1.2cm} & \blank{1.2cm} & \blank{1.2cm} \\
        \hline
        \textbf{Is $V(x) \ge 480$?} & \blank{1.2cm} & \blank{1.2cm} & \blank{1.2cm} \\
        \hline
        \end{tabular}
        \end{center}
        \smallskip

        Solution in interval notation: \blank{5cm}

        Now say it in a sentence, with units.

        \vspace{0.05in}
        \writelines{2}
\end{enumerate}
\end{tcolorbox}

\end{document}
